% Options for packages loaded elsewhere
\PassOptionsToPackage{unicode}{hyperref}
\PassOptionsToPackage{hyphens}{url}
\PassOptionsToPackage{dvipsnames,svgnames,x11names}{xcolor}
\documentclass[
]{article}
\usepackage{xcolor}
\usepackage[margin=1in]{geometry}
\usepackage{amsmath,amssymb}
\setcounter{secnumdepth}{-\maxdimen} % remove section numbering
\usepackage{iftex}
\ifPDFTeX
  \usepackage[T1]{fontenc}
  \usepackage[utf8]{inputenc}
  \usepackage{textcomp} % provide euro and other symbols
\else % if luatex or xetex
  \usepackage{unicode-math} % this also loads fontspec
  \defaultfontfeatures{Scale=MatchLowercase}
  \defaultfontfeatures[\rmfamily]{Ligatures=TeX,Scale=1}
\fi
\usepackage{lmodern}
\ifPDFTeX\else
  % xetex/luatex font selection
  \setmainfont[]{Helvetica Neue}
\fi
% Use upquote if available, for straight quotes in verbatim environments
\IfFileExists{upquote.sty}{\usepackage{upquote}}{}
\IfFileExists{microtype.sty}{% use microtype if available
  \usepackage[]{microtype}
  \UseMicrotypeSet[protrusion]{basicmath} % disable protrusion for tt fonts
}{}
\makeatletter
\@ifundefined{KOMAClassName}{% if non-KOMA class
  \IfFileExists{parskip.sty}{%
    \usepackage{parskip}
  }{% else
    \setlength{\parindent}{0pt}
    \setlength{\parskip}{6pt plus 2pt minus 1pt}}
}{% if KOMA class
  \KOMAoptions{parskip=half}}
\makeatother
\usepackage{longtable,booktabs,array}
\usepackage{caption}
\captionsetup[table]{skip=6pt}
\newcounter{none} % for unnumbered tables
\usepackage{calc} % for calculating minipage widths
% Correct order of tables after \paragraph or \subparagraph
\usepackage{etoolbox}
\makeatletter
\patchcmd\longtable{\par}{\if@noskipsec\mbox{}\fi\par}{}{}
\makeatother
\setlength{\emergencystretch}{3em} % prevent overfull lines
\providecommand{\tightlist}{%
  \setlength{\itemsep}{0pt}\setlength{\parskip}{0pt}}
\usepackage{bookmark}
\IfFileExists{xurl.sty}{\usepackage{xurl}}{} % add URL line breaks if available
\urlstyle{same}
% fallback for those not using the hyperref driver hyperxmp:
\makeatletter
\@ifundefined{xmpquote}{\newcommand{\xmpquote}[1]{#1}}{}
\makeatother
\hypersetup{
  pdftitle={Reporte de Avance No. 1},
  pdfauthor={Equipo de Desarrollo ABAQ},
  colorlinks=true,
  linkcolor={Maroon},
  filecolor={Maroon},
  citecolor={Blue},
  urlcolor={Blue},
  pdfcreator={LaTeX via pandoc}}

\title{Reporte de Avance No. 1}
\author{Equipo de Desarrollo ABAQ}
\date{}

\begin{document}
\maketitle

{
\setcounter{tocdepth}{3}
\tableofcontents
}
\section{Reporte de Avance No.~1}\label{reporte-de-avance-no.-1}

{\def\LTcaptype{none} % do not increment counter
\begin{longtable}[]{@{}ll@{}}
\toprule\noalign{}
\endhead
\bottomrule\noalign{}
\endlastfoot
\textbf{Fecha} & 19 de agosto de 2026 \\
\textbf{Versión} & v1.0.0 \\
\textbf{Estado} & En planeación / Prototipado \\
\textbf{Preparado por} & Equipo de Desarrollo ABAQ \\
\end{longtable}
}

\begin{center}\rule{0.5\linewidth}{0.5pt}\end{center}

\textbf{Descargas} Puedes descargar este documento en los siguientes
formatos:

\begin{verbatim}
-  **[Descargar PDF](latex/report1.pdf)**
-  **[Descargar Word (.docx)](latex/report1.docx)**
\end{verbatim}

\subsection{Resumen ejecutivo}\label{resumen-ejecutivo}

Este documento presenta el primer reporte de avance del proyecto ABAQ.
Durante este período se llevó a cabo la presentación oficial del primer
prototipo funcional de la plataforma a la asociación ABAQ Querétaro,
logrando alinear las expectativas sobre el diseño y la arquitectura del
sistema. Se acordó el desarrollo de un sistema con accesos separados
para administradores y estudiantes, garantizando la seguridad de la
información. Finalmente, se negoció un esquema de trabajo de 8 entregas
quincenales para la liberación del servicio social.

\begin{center}\rule{0.5\linewidth}{0.5pt}\end{center}

\subsection{1. Introducción}\label{introducciuxf3n}

Este documento resume los acuerdos y requerimientos recabados en la
primera reunión oficial de seguimiento (19 de agosto de 2026) entre el
equipo de desarrollo de software y ABAQ Querétaro. El propósito
principal de la sesión fue validar el rumbo técnico del proyecto y
formalizar las condiciones operativas.

\begin{center}\rule{0.5\linewidth}{0.5pt}\end{center}

\subsection{2. Detalles del Acuerdo}\label{detalles-del-acuerdo}

\subsubsection{2.1 Roles por equipo}\label{roles-por-equipo}

El desarrollo del aplicativo está a cargo de un equipo de tres
integrantes. Los tres conforman, de manera conjunta, el \textbf{equipo
de base de datos del aplicativo}, responsable del modelo de datos que
sustenta la plataforma.

{\def\LTcaptype{none} % do not increment counter
\begin{longtable}[]{@{}ll@{}}
\toprule\noalign{}
Integrante & Rol \\
\midrule\noalign{}
\endhead
\bottomrule\noalign{}
\endlastfoot
Fernando Israel Ríos García & Documentación \\
Johan & Developer general \\
Brahn & Developer y diseñador gráfico \\
\end{longtable}
}

El modelo de datos sobre el que trabaja este equipo está implementado en
MongoDB (Mongoose) y se compone, en su estado actual, de las siguientes
colecciones:

\begin{itemize}
\tightlist
\item
  \textbf{Estudiante (\texttt{Student}):} expediente del voluntario ---
  nombre completo, matrícula, correos personal e institucional, escuela,
  carrera, fecha de arranque, auditor asignado, teléfono, fecha de
  nacimiento, horas (a cubrir, acreditadas y pendientes), documentos
  oficiales (carta de aceptación, carta de término, constancia de
  plática), redes sociales, periodo y estado de conclusión.
\item
  \textbf{Contacto (\texttt{Contact}):} seguimiento de comunicaciones
  --- correo de contacto, estudiante asociado, estado y fechas de
  seguimiento.
\item
  \textbf{Encuesta (\texttt{Survey}):} instrumento de bienestar animal
  --- estudiante asociado, datos de mascotas, esterilización,
  vacunación, condiciones de vida y campos de concientización.
\item
  \textbf{Actividad (\texttt{Activity}):} actividades del servicio
  social --- título, descripción, tipo (obligatoria, opcional o
  específica), fechas, valor en horas y horas validadas.
\end{itemize}

\subsubsection{2.2 Fechas de entrega}\label{fechas-de-entrega}

Se acordó un esquema de \textbf{8 sesiones quincenales} (una cada 15
días) a partir de la primera sesión oficial del 19 de agosto de 2026. Se
incluye además la sesión 0 correspondiente a la presesión de
presentación inicial.

{\def\LTcaptype{none} % do not increment counter
\begin{longtable}[]{@{}
  >{\raggedright\arraybackslash}p{(\linewidth - 4\tabcolsep) * \real{0.3333}}
  >{\raggedright\arraybackslash}p{(\linewidth - 4\tabcolsep) * \real{0.3333}}
  >{\raggedright\arraybackslash}p{(\linewidth - 4\tabcolsep) * \real{0.3333}}@{}}
\toprule\noalign{}
\begin{minipage}[b]{\linewidth}\raggedright
Sesión N
\end{minipage} & \begin{minipage}[b]{\linewidth}\raggedright
Fecha
\end{minipage} & \begin{minipage}[b]{\linewidth}\raggedright
Descripción
\end{minipage} \\
\midrule\noalign{}
\endhead
\bottomrule\noalign{}
\endlastfoot
0 & 29 jul 2026 & Presesión de presentación inicial (ver Reporte
No.~0) \\
1 & 19 ago 2026 & Validación del prototipo y definición de
requerimientos (Reporte No.~1) \\
2 & 2 sep 2026 & Avance de seguridad y nuevas pantallas de
administración (documentado en el Reporte No.~2, fechado el 9 sep
2026) \\
3 & 16 sep 2026 & Seguimiento de desarrollo según próximos pasos
acordados \\
4 & 30 sep 2026 & Seguimiento de desarrollo según próximos pasos
acordados \\
5 & 14 oct 2026 & Seguimiento de desarrollo según próximos pasos
acordados \\
6 & 28 oct 2026 & Seguimiento de desarrollo según próximos pasos
acordados \\
7 & 11 nov 2026 & Seguimiento de desarrollo según próximos pasos
acordados \\
8 & 25 nov 2026 & Sesión de cierre y entrega final del sistema \\
\end{longtable}
}

\subsubsection{2.3 Acuerdos obtenidos}\label{acuerdos-obtenidos}

Sobre lo solicitado en la reunión, se obtuvieron los siguientes
acuerdos:

\begin{enumerate}
\def\labelenumi{\arabic{enumi}.}
\tightlist
\item
  \textbf{Entrega del modelo de datos de voluntarios:} se entregará un
  archivo Excel con el modelo de datos de los voluntarios guardados en
  la plataforma.
\item
  \textbf{Cumplimiento de las sesiones quincenales:} se cumplirá con el
  esquema de \textbf{8 sesiones quincenales} hasta diciembre de 2026,
  conforme al conteo iniciado en el primer reporte (19 de agosto de
  2026), culminando con la entrega del sistema completo.
\end{enumerate}

\begin{center}\rule{0.5\linewidth}{0.5pt}\end{center}

\subsection{3. Cambios en el Servidor
(Backend)}\label{cambios-en-el-servidor-backend}

Sin información, no requerido.

\begin{center}\rule{0.5\linewidth}{0.5pt}\end{center}

\subsection{4. Cambios en la Aplicación Web
(Frontend)}\label{cambios-en-la-aplicaciuxf3n-web-frontend}

Sin información, no requerido.

\begin{center}\rule{0.5\linewidth}{0.5pt}\end{center}

\subsection{5. Evidencia Visual}\label{evidencia-visual}

Falta de información.

\begin{center}\rule{0.5\linewidth}{0.5pt}\end{center}

\subsection{6. Estado Actual del
Proyecto}\label{estado-actual-del-proyecto}

\textbf{En Prototipado y Negociación} Se validó el prototipo funcional
inicial y se acordó formalmente la liberación del servicio social a
cambio de la entrega del sistema completo para diciembre de 2026.

\begin{center}\rule{0.5\linewidth}{0.5pt}\end{center}

\subsection{7. Próximos Pasos}\label{pruxf3ximos-pasos}

\begin{enumerate}
\def\labelenumi{\arabic{enumi}.}
\tightlist
\item
  \textbf{Vistas Separadas:} Desarrollar y presentar el panel dividido
  entre Administrador y Estudiante.
\item
  \textbf{Datos de Prueba:} Alimentar la plataforma con registros de
  prueba realistas.
\item
  \textbf{Documentación Formal:} Elaborar un documento formal que
  detalle la propuesta técnica y el calendario de entregas.
\end{enumerate}

\end{document}
